\chapter{Objetivos}

\section{Objetivo Geral}

Formar profissionais capazes de projetar, desenvolver, integrar, implantar e manter soluções baseadas em Inteligência Artificial, Ciência de Dados, Machine Learning, Deep Learning, Inteligência Artificial Generativa, Agentes Inteligentes, Internet das Coisas (IoT), Computação Distribuída e MLOps, aplicando metodologias modernas de Engenharia de Software e Ciência de Dados para resolução de problemas complexos em ambientes corporativos, industriais e tecnológicos.

Ao término da formação, o egresso deverá ser capaz de atuar em todas as etapas do ciclo de vida de soluções inteligentes, desde a identificação do problema até a implantação, monitoramento, manutenção e evolução contínua dos modelos e aplicações desenvolvidas.

\section{Objetivos Específicos}
Ao concluir o curso, o estudante deverá ser capaz de:

\textbf{Fundamentos}
\begin{itemize}
\item Compreender os princípios científicos da Inteligência Artificial;
\item Aplicar fundamentos matemáticos utilizados em Machine Learning;
\item Utilizar programação em Python para desenvolvimento de soluções inteligentes;
\item Empregar ferramentas modernas de desenvolvimento colaborativo.
\end{itemize}

\textbf{Ciência de Dados}
\begin{itemize}
\item Coletar dados provenientes de diferentes fontes;
\item Preparar bases de dados para treinamento de modelos;
\item Aplicar técnicas de limpeza e transformação de dados;
\item Desenvolver processos de Engenharia de Dados;
\item Construir pipelines de dados.
\end{itemize}

\textbf{Machine Learning}
\begin{itemize}
    \item Desenvolver modelos supervisionados.
    \item Desenvolver modelos não supervisionados.
    \item Selecionar algoritmos adequados para diferentes problemas.
    \item Avaliar desempenho de modelos.
    \item Interpretar métricas estatísticas.
    \item Aplicar técnicas de Explainable AI.
\end{itemize}

\textbf{Deep Learning}
\begin{itemize}
    \item Desenvolver redes neurais artificiais.
    \item Utilizar TensorFlow e PyTorch.
    \item Aplicar CNNs.
    \item Aplicar Transformers.
    \item Desenvolver modelos multimodais.
    \item Executar Fine-Tuning de modelos.
\end{itemize}

\textbf{Inteligência Artificial Generativa}
\begin{itemize}
    \item Utilizar Large Language Models.
    \item Desenvolver aplicações utilizando APIs de IA.
    \item Criar copilotos inteligentes.
    \item Construir aplicações baseadas em IA Generativa.
    \item Desenvolver soluções multimodais.
\end{itemize}

\textbf{Engenharia de Prompt}
\begin{itemize}
    \item Construir prompts eficientes.
    \item Desenvolver Engenharia de Contexto.
    \item Aplicar padrões avançados de prompting.
    \item Desenvolver estratégias para aumento da qualidade das respostas produzidas por LLMs.
\end{itemize}

\textbf{Agentes Inteligentes}
\begin{itemize}
    \item Projetar arquiteturas orientadas a agentes.
    \item Desenvolver agentes inteligentes.
    \item Construir sistemas RAG.
    \item Desenvolver agentes utilizando LangGraph.
    \item Desenvolver aplicações utilizando MCP.
    \item Integrar agentes a APIs corporativas.
    \item Desenvolver sistemas multiagentes.
\end{itemize}

\textbf{Desenvolvimento de Software}
\begin{itemize}
    \item Utilizar IA Generativa durante todo o ciclo de desenvolvimento.
    \item Automatizar testes.
    \item Automatizar documentação.
    \item Desenvolver APIs.
    \item Desenvolver aplicações Web.
    \item Desenvolver aplicações Full Stack.
    \item Automatizar pipelines de integração contínua.
\end{itemize}

\textbf{Visão Computacional}
\begin{itemize}
    \item Desenvolver aplicações utilizando OpenCV.
    \item Desenvolver modelos YOLO.
    \item Desenvolver sistemas OCR.
    \item Desenvolver soluções para classificação e detecção de objetos.
\end{itemize}

\textbf{IoT e Computação Distribuída}
\begin{itemize}
    \item Integrar dispositivos inteligentes.
    \item Desenvolver soluções utilizando MQTT.
    \item Utilizar OPC-UA.
    \item Desenvolver arquiteturas Edge Computing.
    \item Utilizar Kafka.
    \item Utilizar Spark.
    \item Utilizar Kubernetes.
    \item Desenvolver aplicações distribuídas.
\end{itemize}

\textbf{MLOps}
\begin{itemize}
    \item Implantar modelos em produção.
    \item Automatizar pipelines de treinamento.
    \item Implementar monitoramento.
    \item Gerenciar versionamento de modelos.
    \item Detectar Data Drift.
    \item Aplicar governança de IA.
\end{itemize}

\textbf{Competências Profissionais}
\begin{itemize}
    \item Trabalhar em equipes multidisciplinares.
    \item Comunicar resultados técnicos.
    \item Desenvolver pensamento crítico.
    \item Aplicar princípios éticos na utilização da Inteligência Artificial.
    \item Desenvolver soluções alinhadas às necessidades das organizações.
\end{itemize}